\documentclass[12pt]{article}  % Tipo de archivo
\usepackage[utf8]{inputenc}    % Codificación utf-8
\usepackage[spanish]{babel}    % Idioma español
\usepackage{graphicx}          % Necesario para imágenes

\title{
    \textsc{\Large Universidad Nacional Autónoma de México\\[0.2cm]
        {\normalsize Facultad de Ciencias, 2027-I}\\[0.2cm]
        {\normalsize Modelado y Programación}}\\[0.6cm]
    \rule{\linewidth}{0.4pt}\\[0.4cm]
    {\textsc{\Large {\color{maincolor}Proyecto:} Base de Datos Musical}}\\[0.1cm]
    \rule{\linewidth}{0.4pt}\\[1em]
}
\author{
    \textsc{Alumno:} Valeria Becerril Acosta\\[0.2cm]
    \textsc{Número de cuenta:} 323128489\\[0.4cm]
}
\date{\today\\[2.6cm]}

% PAQUETES DE ESTILO
\usepackage{xcolor}           % Para definir y usar colores
\usepackage{fancyhdr}         % Para los encabezados y pies de página fancy
\usepackage{titlesec}         % Para configurar \titleformat
\usepackage{enumitem}         % Para configurar \setlist
\usepackage{lmodern}          % Para la familia de fuentes
\usepackage[T1]{fontenc}      % Para codificación de fuente
\usepackage{listings}         % Para el estilo de código mystyle
\usepackage[most]{tcolorbox}  % Para los cuadros de color
\usepackage{tabularx}         % Para el ancho de tablas
\usepackage{hyperref}         % Para los enlaces

% CONFIGURACIONES Y COLORES
\definecolor{maincolor}{RGB}{0, 51, 102}
\definecolor{codegreen}{rgb}{0,0.6,0}
\definecolor{codegray}{rgb}{0.5,0.5,0.5}
\definecolor{codepurple}{rgb}{0.58,0,0.82}
\definecolor{foldercolor}{RGB}{255, 193, 7}
\definecolor{pdfcolor}{RGB}{220, 53, 69}
\definecolor{readmecolor}{RGB}{40, 167, 69}
\definecolor{txtcolor}{RGB}{0, 123, 255}

% TARJETA CRC
\newcommand{\crc}[3]{
    \begin{center}
        \begin{tcolorbox}[
            colback=maincolor!10, colframe=maincolor,
            colupper=black, fontupper=\footnotesize\ttfamily,
            boxrule=0.2mm, arc=2pt,
            left=3pt, right=3pt, top=1pt, bottom=1pt,
            title={\fontseries{b}\selectfont\footnotesize Class
                \begin{center}
                    #1
                \end{center}},
            colbacktitle=maincolor!35,
            coltitle=black,
            toptitle=4pt, bottomtitle=4pt
            ]

            \begin{tabularx}{\linewidth}{@{}X|X@{}}
              \textbf{\footnotesize Responsabilidades}
              & \textbf{\footnotesize Colaboradores} \\[4pt]
              \begin{itemize}[itemsep=-5pt, leftmargin=*]
                  #2
              \end{itemize}
              &  \begin{itemize}[itemsep=-5pt, leftmargin=*]
                  #3
              \end{itemize}
            \end{tabularx}
        \end{tcolorbox}
    \end{center}
}

% ESTILO DE CÓDIGO
\lstdefinestyle{mystyle}{
    backgroundcolor=\color{maincolor!10},
    commentstyle=\color{codegreen},
    keywordstyle=\color{magenta},
    numberstyle=\tiny\color{codegray},
    stringstyle=\color{codepurple},
    basicstyle=\ttfamily\footnotesize,
    breakatwhitespace=false,
    breaklines=true,
    captionpos=b,
    keepspaces=true,
    numbers=left,
    numbersep=5pt,
    showspaces=false,
    showstringspaces=false,
    showtabs=false,
    tabsize=2
}
\lstset{style=mystyle}

% CONFIGURACIÓN DE ENCABEZADO Y PIE DE PÁGINA
\pagestyle{fancy}
\fancyhead[LO,L]{UNAM}
\fancyhead[RO,R]{\textit{\textsc{Proyecto}}}
\fancyfoot[LO,L]{Facultad De Ciencias, 2027-I}
\fancyfoot[CO,C]{\thepage}
\fancyfoot[RO,R]{Modelado y Programación}
\renewcommand{\headrulewidth}{0.4pt}
\renewcommand{\footrulewidth}{0.4pt}

% CAMBIAR COLOR DE SECCIONES
\titleformat{\section}
{\color{maincolor}\normalfont\Large\bfseries}
{\thesection}{1em}{}

\titleformat{\subsection}
{\color{maincolor}\normalfont\large\bfseries}
{\thesubsection}{1em}{}

\titleformat{\subsubsection}
{\color{maincolor}\normalfont\normalsize\bfseries}
{\thesubsubsection}{1em}{}

% CAMBIAR COLOR DE ÍTEMS EN LISTAS
\setlist[enumerate]{label=\textcolor{maincolor}{\arabic*.}, left=1.5em}
\setlist[itemize]{label=\textcolor{maincolor}{\alph*}, left=1.5em}
\setlist[itemize,1]{label=\textcolor{maincolor}{\textbullet}, left=1.5em}


% INICIO DEL DOCUMENTO
\begin{document}

\maketitle
\thispagestyle{empty}

\begin{abstract}
    Este proyecto presenta una base de datos musical en C++17 gestionada con
    SQLite y diseñada bajo Qt, usando el patrón de diseño MVC y DAO, apoyándose de
    herramientas como Meson para la construcción, Doxygen para la documentación y
    bibliotecas como TagLib para la lectura de los metadatos MP3.
\end{abstract}

\newpage
\tableofcontents

\newpage
\section{Introducción}

Este proyecto busca implementar una base de datos musical en C++ que guarde la
información de las canciones de un directorio, así como sus artistas, albumes y
sus relaciones entre sí de manera local. Para ello, se usa SQLite para gestionar
la base de datos y Qt con QML para desarrollar la interfaz.

\section{Problema}

Esencialmente, el desarrollo de este proyecto se reduce a los siguientes 4
componentes principales:

\begin{itemize}
  \item {\color{maincolor}\textbf{El Minero}}\\
    Se encarga de recorrer el directorio elegido para identificar y leer los
    datos de los archivos .mp3 que este contiene.
  \item {\color{maincolor}\textbf{La Interfaz}}\\
    La vista que permite al usuario ver y modificar la información de la base de
    datos.
  \item {\color{maincolor}\textbf{La Base de Datos}}\\
    El motor que se encarga de guardar toda la información extraída de los
    archivos .mp3 y los cambios realizados por el usuario.
  \item {\color{maincolor}\textbf{Extra: El Mini Compilador}}\\
    Un intérprete de comandos capaz de permitir búsquedas avanzadas.
\end{itemize}

\section{Análisis}

\subsection{El Minero}

Recorre un directorio arbitrario para identificar y leer los datos de los
archivos .mp3 que este contiene.

Para ello, primero se obtiene toda la información que ya se tiene localmente en
\texttt{base\_datos.db} (canciones, álbumes y artistas). Luego, se compara con
la nueva información extraida del directorio y se almacena la faltante en
objetos que los DAOs corresppndientes recibiran para agregarlos a la base de datos.

\subsection{La Interfaz}

Muestra la información de la base de datos y permite modificarla sin alterar los
archivos originales en el directorio.

Cuenta con botones para minar el directorio, editar y reproducir una canción,
así como deslizadores para cambiar el volumen y la posición de la canción.

Incluye un pequeño formulario de buscar que permite filtar por cada campo
mostrado de una canción: título, artista, álbum, año y género. Seleccionar una
canción permite eliminarla o abrir un formulario donde editarla, según el
botón seleccionado.

\subsection{La Base de Datos}

El motor encargado de guardar toda la información extraída de los .mp3 en
el archivo \texttt{base\_datos.db}. Esta se puede modificar de acuerdo a los cambios
realizados por el usuario desde la interfaz.

Para mantenerla se tiene un objeto que unicamente gestiona la conexión y para
modificarla se usan objetos que representaran a las canciones, álbumes y
artistas, así como sus respectivos DAOs para obtener y modificar la información
de la base de datos.

\subsection{El Mini Compilador}

Se omitio por falta de tiempo y en su lugar se implementó el formulario de búsqueda
mencionado de la interfaz.

\section{Diseño}

El programa inicia su ejecución con un método main que solo se encarga de
atrapar las expeciones, mostrar los mensajes de error e iniciar la Aplicación.

\crc{Aplicacion}{
  \item Válida los argumentos de entrada del programa.
  \item Carga el controlador y la interfáz gráfica.
}{
  \item Controlador
}

El \texttt{Controlador} es el objeto que se encarga de gestionar las acciones
principales del programa. Recibe los datos de la vista y manda a llamar a los
métodos necesarios para llevar a cabo las acciones solicitadas por el usuario.

\crc{Controlador}{
  \item Inicia la base de datos.
  \item Maneja los DAOs.
  \item Recibe y envia información a la interfaz.
}{
  \item «namespace» Minero
  \item BaseDeDatos
  \item RolaDao
  \item AlbumDao
  \item ArtistaDao
  \item PersonaDao
}

El espacio de nombres \texttt{Minero} tiene una única función para minar el
directorio elegido por el usuario. Su trabajo es recorrer el directorio y
comparar la información obtenida de el con la que ya pertenece a la base de
datos para obtener todas las canciones, álbumes y artistas faltantes. Estos
objetos se guardan en la estructura \texttt{Elementos} y son manejados por los
DAOs correspondientes.

La clase \texttt{BaseDeDatos} tiene una única función: manejar la conexión de la
base de datos. De esta forma, el \texttt{Controlador} nunca conoce la conexión y
los DAOs solo interactuan con la conexión al realizar consultas.

\crc{BaseDeDatos}{
  \item Crea el archivo \texttt{base\_datos.db} si no existe.
  \item Maneja la conexión de la base de datos.
}{
  \item N/A
}

El DAO de las canciones realiza todas las consultas SQL referentes a las
canciones. Es decir, se encarga de añadir, obtener, actualizar, eliminar y
buscar canciones.

\crc{RolaDao}{
  \item Realiza las consultas SQL de las canciones.
  \item Obtiene y modifica los datos de las canciones en la base de datos.
}{
  \item Rola
  \item «struct» RolaNueva
  \item «struct» VistaRola
  \item «struct» ParamRola
  \item «namespace» DaoUtils
}

Las estructuras colaboradoras de \texttt{RolaDao} son objetos que guardan datos
relacionados con la información de las canciones.

Una \texttt{Rola} guarda todos los datos de una canción perteneciente a la tabla
SQL \texttt{rolas}. Una \texttt{RolaNueva} envuelve una \texttt{Rola} y guarda
la ruta del álbum y el nombre del artista, información extra requerida
para relacionar una canción con otras tablas SQL. De manera similar, una
\texttt{VistaRola} guarda adicionalemente el nombre del album y del artista como
la información que se muestra al usuario.

Por otro lado, una \texttt{ParamRola} guarda todos los parámetros (opcionales) por
los que se puede buscar una canción: título, año, genéro, álbum y artista.

Continuando con los DAOs faltantes, \texttt{AlbumDao}, \texttt{ArtistaDao} y
\texttt{PersonaDao}, como las responsabilidades que tienen son similares, nos
limitaremos a mostrar la tarjeta CRC del DAO de los álbumes.

\crc{AlbumDao}{
  \item Realiza las consultas SQL de los álbumes.
  \item Obtiene y modifica los datos de los álbumes en la base de datos.
}{
  \item Album
  \item «namespace» DaoUtils
}

De manera similar a una \texttt{Rola}, un \texttt{Album} guarda los datos de un
album en la tabla SQL \texttt{albums}, un \texttt{Artista} guarda los datos de
un artista en la tabla SQL \texttt{performers} y una \texttt{Persona} guarda los
datos de un artista en la tabla SQL \texttt{persons}.

Todos los DAOs se apoyan del espacio de nombres \texttt{DaoUtils} que tiene
funciones auxiliares para las consultas SQL.

Por su parte, la interfaz se compone de cuatro objetos principales: la ventana
principal (\texttt{Main}) donde se muestra la aplicación, el encabezado
(\texttt{HeaderBar}) con los botones principales, el pie (\texttt{FooterBar})
con los controles para reproducir una canción y el cuerpo de la ventana
(\texttt{Body}) que muestra la información de las canciones en la base de
datos. Estos componentes interactuan con los método expuestos en el
\texttt{Controlador}.

\section{Implementación}

Para la implementación de este proyecto se usó Meson (con Ninja) como sistema de
construcción, creando varios archivos \texttt{meson.build}, uno en cada
directorio principal para gestionar los archivos, las dependencias y los
recursos.

Entre las dependencias se encuentran Qt6 y SQLite3, esenciales para la
compilación y ejecución del programa. Asimismo, se usó Doxygen para generar la
documentación, pero se puede prescindir de esta herramienta sin que el
funcionamiento del programa se vea afectado.

Primero, compilamos:

\begin{center}
    {\texttt{meson setup build \\
            meson compile -C build}}
\end{center}

Luego, para correr el programa:

\vspace{1.2em}
\centerline{\texttt{./build/src/main/db/musical-db ruta/archivos/mp3}}
\vspace{1.2em}

Donde \texttt{ruta/archivos/mp3} es el directorio que el programa minará.

\section{Conclusión}

El desarrollo de este proyecto permitió integrar distintos componentes para
construir una aplicación capaz de gestionar una colección de música de manera
local. La separación de responsabilidades entre la interfaz, el controlador, el
minero y los DAOs facilitó la organización y desarrollo del programa.

Aunque el mini compilador no se implementó, se consiguió desarrollar una
aplicación funcional que cumple con los objetivos principales del proyecto.

% FIN DEL DOCUMENTO
\end{document}
